% 21_body.tex — 산출물 ㉑ 창업 헌장·공동창업자 합의서 (빈 템플릿 / 완성 예시 공용 본문) — gen_product_templates.py 가 Product_specs.py 에서 생성
% 드라이버: 21_창업헌장합의서_빈템플릿.tex (\wbblanktrue) / 21_창업헌장합의서_완성예시.tex (\wbblankfalse — 워크북 §X.5 확정 후)
\documentclass[10pt,a4paper]{article}
\usepackage[a4paper,margin=16mm,top=14mm,bottom=20mm]{geometry}
\def\DocNum{21}\def\DocDay{5}\def\DocIdx{1}
\def\DocTitle{창업 헌장·공동창업자 합의서}
\def\DocSub{Founders' Charter \& Co-founder Agreement}
\def\DocVariant{\B{빈 템플릿}{딥게이지 완성 예시}}
\usepackage{wbstyle}
\def\DocCirc{㉑}   % newunicodechar(㉛~㊵ 폴백) 뒤에 정의해야 활성 문자로 읽힌다
\begin{document}
\docheader
 {\Bg{[회사명 · 제품명]}{주식회사 딥게이지 · GAUGE(가칭)}}
 {\Bg{[v0.x / D+n]}{v1.0 / 2026-03-13 (D+12)}}
 {\Bg{[작성자 3인]}{강민수 · 오세진 · 윤하경\rd{7mm}}}
 {\Bg{[검토자 — 법무·투자자]}{(법무 검토 예정 · 노들투자파트너스)\rd{7mm}}}

%% ------------------------------------------------------------------ 1
\secbar{1. 문서 정보}
\guide{회사·설립일·버전·서명자를 적는다. 합의서는 시드 텀시트(㉚) 전에 확정되어야 한다 — 투자자가 먼저 묻는 것이 이 문서다.}
{\footnotesize
\begin{tabularx}{\linewidth}{|L{34mm}|Y|}\hline
\hd{항목} & \hd{내용} \\ \hline
\ifwbblank
\rd{8mm}회사명 / 설립일 &  \\ \hline
\rd{8mm}제품명 / 한 줄 정의 &  \\ \hline
\rd{8mm}문서 버전 / 기준일 &  \\ \hline
\rd{8mm}공동창업자(직위) &  \\ \hline
\rd{8mm}다음 갱신 시점 &  \\ \hline
\else
회사명 / 설립일 & 주식회사 딥게이지(DeepGauge Inc.) / 2026-03-02, 자본금 1억(액면 100원 × 1,000,000주) \\ \hline
제품명 / 한 줄 정의 & GAUGE(가칭) — 검사원이 '라인이 도는 동안 나중에 책임을 물을 수 없게 기록을 남기기 위해' 고용하는, 사진·음성에서 검사기록과 부적합 처리를 만들어 주는 서비스 \\ \hline
문서 버전 / 기준일 & v1.0 / 2026-03-13 \\ \hline
공동창업자(직위) & 강민수(대표이사) · 오세진(CTO) · 윤하경(CPO) \\ \hline
다음 갱신 시점 & 시드 텀시트 수령 시(2026-08 예정, D+150) — 항목 3·4를 텀시트와 대조해 v1.1 \\ \hline
\fi
\end{tabularx}}

%% ------------------------------------------------------------------ 2
\secbar{2. 왜 이 회사인가 — 문제와 팀의 자격}
\guide{해법이 아니라 문제를 적는다. 숫자 하나 이상(작성 47분·전기 31분·통보 3.7일·8D 11.5시간·클레임 2,700만). 팀이 이 문제를 풀 자격(현장 경험·기술·구축 경험)을 한 문단. 온담 이관 목록 6번과의 관계는 33개월 뒤의 일이다 — 지금은 쓰지 않는다.}
\begin{fieldbox}
\ifwbblank
\lines{8}{9.5mm}
\else
자동차 2·3차 협력사의 검사원은 하루 47분을 검사기록을 쓰는 데 쓰고, 그중 31분은 라인에서 이미 적은 값을 종이 → 엑셀 → 1차사 포털로 다시 적는 시간이다(G-33). 부적합이 발생하면 사진과 근거를 모으는 데 이틀이 걸려 통보가 평균 3.7일 늦고(G-34), 8D 한 건에 11.5시간을 쓴다(G-35). 클레임 한 건은 평균 2,700만 원이다(G-36). 공장 한 곳의 전기·중복 입력만 월 약 126만 원어치 사람 시간이다.

지금인 이유는 둘 — 사진 한 장에서 치수·외관 결함을 읽는 비전-언어 모델이 실용 정확도에 도달했고 그것을 만들어 본 사람이 이 팀에 있다; 1차사들이 2·3차의 검사기록을 포털로 제출하게 요구하기 시작했다.

팀의 자격: 강민수(완성차 QA 9년 + 2차사 품질팀장 3년 — 네 숫자를 잰 사람), 오세진(포털 비전팀 5년, VLM·OCR 배포), 윤하경(스마트팩토리 SI PM 4년). 우리가 모르는 것 — 검사원이 이것을 원하는가. 27건 인터뷰가 답한다.
\fi
\end{fieldbox}

%% ------------------------------------------------------------------ 3
\secbar{3. 초기 지분·주식수·옵션풀}
\guide{주식수와 지분을 함께 적고, 옵션풀이 pre-money인지 post-money인지를 반드시 적는다 — 미명시가 이 문서의 첫 함정이다.}
{\footnotesize
\begin{tabularx}{\linewidth}{|L{30mm}|L{26mm}|L{22mm}|Y|}\hline
\hd{주주} & \hd{주식수} & \hd{지분(FD)} & \hd{출자·역할 근거} \\ \hline
\ifwbblank
\rd{8mm}강민수 CEO &  &  &  \\ \hline
\rd{8mm}오세진 CTO &  &  &  \\ \hline
\rd{8mm}윤하경 CPO &  &  &  \\ \hline
\rd{8mm}옵션풀(미발행 유보) — pre / post: \_\_\_\_\_\_ &  &  &  \\ \hline
\rd{8mm}계 &  &  &  \\ \hline
\else
강민수 CEO & 450,000 & 45.0\% & 풀타임 · 개인 출자 4,000만 · 문제의 원 소유자 · 고객 관계 \\ \hline
오세진 CTO & 300,000 & 30.0\% & 풀타임 · 출자 1,500만 · 핵심 기술(모델) \\ \hline
윤하경 CPO & 150,000 & 15.0\% & 풀타임 · 출자 500만 · 발견·범위·구축 경험 · D+0 합류 가장 늦음 \\ \hline
옵션풀 — 설립 시 FD 포함(pre) · 시드 추가 풀은 post-money 요구 & 100,000 & 10.0\% & 첫 5명 채용용. pre 기준 추가 풀이면 그 희석은 세 사람 몫 — 협상 요청 1번 \\ \hline
계 & 1,000,000 & 100.0\% & K-01 \\ \hline
\fi
\end{tabularx}}

%% ------------------------------------------------------------------ 4
\secbar{4. 베스팅·클리프·가속}
\guide{기간·클리프·월 단위 베스팅·가속 조건(인수·해임)을 창업자마다 적는다. 정본은 미확정 상태로 시작한다 — 여러분이 정한다.}
{\footnotesize
\begin{tabularx}{\linewidth}{|L{26mm}|L{18mm}|L{18mm}|L{22mm}|Y|}\hline
\hd{창업자} & \hd{총 기간} & \hd{클리프} & \hd{베스팅 단위} & \hd{가속 조건(single/double trigger)} \\ \hline
\ifwbblank
\rd{9mm}강민수 &  &  &  &  \\ \hline
\rd{9mm}오세진 &  &  &  &  \\ \hline
\rd{9mm}윤하경 &  &  &  &  \\ \hline
\else
강민수 & 4년 & 1년(2027-03-02) & 월 1/48 & double trigger 100\% · single 없음 \\ \hline
오세진 & 4년 & 1년 & 월 1/48 & 동일 \\ \hline
윤하경 & 4년 & 1년 & 월 1/48 & 동일 \\ \hline
\fi
\end{tabularx}}

%% ------------------------------------------------------------------ 5
\secbar{5. 의사결정 규칙}
\guide{만장일치 영역 / 과반 영역 / 대표 전결 영역을 구분하고 교착 시 해소 절차(외부 조언자·투자자·동전)를 적는다.}
{\footnotesize
\begin{tabularx}{\linewidth}{|L{50mm}|L{34mm}|Y|}\hline
\hd{결정의 종류} & \hd{규칙(만장일치·과반·전결)} & \hd{교착 해소} \\ \hline
\ifwbblank
\rd{8mm} &  &  \\ \hline
\rd{8mm} &  &  \\ \hline
\rd{8mm} &  &  \\ \hline
\rd{8mm} &  &  \\ \hline
\rd{8mm} &  &  \\ \hline
\rd{8mm} &  &  \\ \hline
\rd{8mm} &  &  \\ \hline
\else
지분 변동·투자·합병·해산·창업자 보상 & 만장일치 & 결렬 시 보류(진행 안 함이 기본값) \\ \hline
채용·해고(창업자 외) & 과반(2/3) & — \\ \hline
범위 절단·Won't·가격 & 과반 + CPO 거부권 1회/분기(범위) · CTO 거부권 1회/분기(아키텍처) & 거부권 시 2주 내 재논의, 재논의는 과반 \\ \hline
500만 이하 지출·가격표 안 계약 & 대표 전결, 월 1회 사후 보고 & — \\ \hline
500만 초과·가격표 밖(커스텀·온프렘) & 과반 & — \\ \hline
어디에도 없는 것 & 대표 결정, 주간 회의 보고 & 두 사람 반대 시 만장일치 항목으로 승격 \\ \hline
교착(2:1 두 번 반복) & 멘토 1인 서면 의견 후 대표 결정 & 6개월 뒤 재검토 \\ \hline
\fi
\end{tabularx}}

%% ------------------------------------------------------------------ 6
\secbar{6. 역할과 책임 — 그리고 용어 합의}
\guide{CEO/CTO/CPO가 각각 무엇을 최종 결정하는가. 이 트랙에서 PM은 프로덕트 매니저다 — 조가 쓸 용어(PM/PO/PdM)를 한 줄로 합의해 적는다.}
{\footnotesize
\begin{tabularx}{\linewidth}{|L{22mm}|Y|L{50mm}|}\hline
\hd{직위} & \hd{최종 결정하는 것} & \hd{하지 않는 것} \\ \hline
\ifwbblank
\rd{9mm}CEO 강민수 &  &  \\ \hline
\rd{9mm}CTO 오세진 &  &  \\ \hline
\rd{9mm}CPO 윤하경 &  &  \\ \hline
\rd{9mm}용어 합의: PM = \_\_\_\_\_\_\_\_ &  &  \\ \hline
\else
CEO 강민수 & 고객(어느 공장에 먼저) · 자본 · 채용 최종 · 대외 & 코드·모델 리뷰 · 화면 설계 \\ \hline
CTO 오세진 & 아키텍처 · 모델·데이터셋·eval · 인프라·보안 & 가격 · 고객 약속(납기·기능) \\ \hline
CPO 윤하경 & 발견(인터뷰·트리·가정) · 범위 · 가격 제안 · 로드맵 & 채용 최종 · 아키텍처 \\ \hline
용어 합의: PM = 프로덕트 매니저(CPO 조직) & PO 직함은 쓰지 않는다 · 프로젝트 관리자는 '프로젝트 PM' & 첫 PM 채용(D+300) 시 CPO와의 경계를 이 표에 추가 \\ \hline
\fi
\end{tabularx}}

%% ------------------------------------------------------------------ 7
\secbar{7. 이탈·해지 처리}
\guide{good leaver / bad leaver 정의, 미베스팅 주식 처리(소각·회수 가격), 경업·비밀유지 기간. '그런 일은 없을 것'이라고 적는 조가 D+230에 첫 이탈을 만난다.}
{\footnotesize
\begin{tabularx}{\linewidth}{|L{40mm}|Y|L{44mm}|}\hline
\hd{상황} & \hd{주식 처리} & \hd{기타 조건} \\ \hline
\ifwbblank
\rd{9mm}자발 퇴사(good leaver) &  &  \\ \hline
\rd{9mm}해임(bad leaver) &  &  \\ \hline
\rd{9mm}사망·질병 &  &  \\ \hline
\rd{9mm}경업·비밀유지 &  &  \\ \hline
\else
자발 퇴사(good leaver) & 미베스팅 소각 · 베스팅분 유지, 회사 우선매수권(최근 라운드 주당가, 라운드 전이면 액면가 10배) & 인수인계 1개월 \\ \hline
해임(bad leaver) & 미베스팅 소각 + 베스팅분 액면가 회수 & 당사자 제외 2인 만장일치 \\ \hline
사망·질병(6개월 이상) & 100\% 가속, 상속·본인 보유 & — \\ \hline
경업·비밀유지 & 경업 1년 · 비밀유지 3년 & 위반 시 bad leaver \\ \hline
\fi
\end{tabularx}}

%% ------------------------------------------------------------------ 8
\secbar{8. 자금 계획 요약}
\guide{현금·월 고정비·런웨이·다음 마일스톤(무엇이 되면 다음 자금을 받을 수 있는가)을 적는다. 정본 §3.1.}
{\footnotesize
\begin{tabularx}{\linewidth}{|L{40mm}|L{34mm}|Y|}\hline
\hd{항목} & \hd{값} & \hd{근거·출처} \\ \hline
\ifwbblank
\rd{8mm}현금(설립 시) &  &  \\ \hline
\rd{8mm}월 고정비 &  &  \\ \hline
\rd{8mm}런웨이(개월) &  &  \\ \hline
\rd{8mm}다음 마일스톤 &  &  \\ \hline
\rd{8mm}마일스톤까지 필요한 것 &  &  \\ \hline
\else
현금(설립 시) & 1억 6,000만 & 개인 출자 6,000(4,000/1,500/500) + 예비창업패키지 5,000 + 자본금 잔여 5,000 (G-15) \\ \hline
월 고정비 & 2,300만(3명) → 3,230만(5명, D+70~) & 급여 1,050 + 코워킹 180 + 클라우드·API 420 + 기타 650 [추가 설정]; 채용 +550 +380 \\ \hline
런웨이(개월) & 7.0(3명) → D+90 잔액 약 8,600만, 5명 기준 2.7 & day5\_discovery.py \\ \hline
다음 마일스톤 & D+90 LOI 2건 + WoZ·OCR 벤치 결과 → 노들 선투자 2억(D+150) → 시드 12억(D+180) & 정본 §3.2 \\ \hline
마일스톤까지 필요한 것 & 인터뷰 27건 · 파일럿 공장 2곳 협조 · 도면 300장 & 항목 2 '우리가 모르는 것' \\ \hline
\fi
\end{tabularx}}

%% ------------------------------------------------------------------ 9
\secbar{9. 서명 — 그리고 역할 전환 선언}
\guide{'당신은 어제까지 이 계약의 발주자였고, 오늘부터 공급자다.' 조원 전원이 서명한다.}
{\footnotesize
\begin{tabularx}{\linewidth}{|Y|Y|Y|Y|}\hline
\hd{강민수 CEO} & \hd{오세진 CTO} & \hd{윤하경 CPO} & \hd{(조원)} \\ \hline
\ifwbblank
\rd{16mm} & & &  \\ \hline
\else
\rd{16mm} 강민수 (서명) & 오세진 (서명) & 윤하경 (서명) & (조원 서명) \\ \hline
\fi
\rd{7mm}일자 & & &  \\ \hline
\end{tabularx}}

\end{document}